\documentclass[UTF8]{ctexart}
\usepackage{geometry}
\geometry{a4paper, margin=2.5cm}
\usepackage{listings}
\usepackage{xcolor}
\usepackage{fontspec}
\setmonofont{DejaVu Sans Mono}  % 若系统无此字体，可改为 Courier New 或 SimHei
\lstset{
    basicstyle=\ttfamily\small,
    breaklines=true,
    frame=single,
    numbers=left,
    numberstyle=\tiny,
    tabsize=4,
    language=C,
    keywordstyle=\color{blue},
    commentstyle=\color{green!60!black},
    stringstyle=\color{red}
}

\title{代码文件括号合法性批量校验工具实验报告}
\author{}

\begin{document}
\maketitle

\begin{center}
    \textbf{班级：}2024217802\quad \textbf{姓名：} 严子竣\quad \textbf{学号：}2024212622
\end{center}

\section{需求分析}

\subsection{程序功能}
本程序模拟软件开发工程实践中静态代码分析工具的核心逻辑——对代码文件中的括号进行合法性校验。要求能够批量处理指定目录下的 C/C++/Java 源文件（扩展名为 .c、.cpp、.java），检验其中的圆括号 ‘()’、方括号 ‘[]’、花括号 ‘{}’是否匹配正确。

程序需要满足以下具体功能：
\begin{itemize}
    \item 读取用户输入的目录路径，自动遍历该目录下的所有目标文件。
    \item 对每个文件逐行扫描，忽略注释（实验假设文件中没有注释），使用栈结构进行括号匹配。
    \item 若文件括号完全匹配（每个左括号有对应的右括号且嵌套顺序正确），输出“文件[文件名]括号匹配合法”。
    \item 若发现非法情况，输出首个错误所在的行号及具体错误类型，包括：
        \begin{itemize}
            \item 未匹配的右括号（多余右括号）
            \item 括号嵌套顺序错误（如 ‘([)]’）
            \item 未匹配的左括号（文件结束后栈中仍有左括号）
            \item 括号嵌套过深（超过预设最大深度 10000 层）
        \end{itemize}
    \item 每个文件独立校验，一个文件的错误不影响后续文件的检查。
\end{itemize}

\subsection{输入形式与值的范围}
\begin{itemize}
    \item 用户通过键盘输入一个合法的 Windows 目录路径（如 ‘D:/test’ 或 ‘C:/project/src‘）。
    \item 程序仅处理该目录下（不递归子目录）的 ‘.c’、‘.cpp’、‘.java’ 文件。
    \item 每个源文件的行长度不超过 4096 字符，括号嵌套深度不超过 10000 层。
\end{itemize}

\subsection{输出形式}
\begin{itemize}
    \item 合法时：\texttt{文件[完整路径]括号匹配合法}
    \item 非法时：\texttt{第X行：错误类型描述}（错误类型包括“未匹配的右圆括号”“方括号嵌套顺序错误”“未匹配的左花括号”等）
    \item 若目录不存在或无目标文件，输出相应提示信息。
\end{itemize}

\subsection{测试数据}
\begin{itemize}
    \item \textbf{正确输入示例}：路径 ‘D:/test’，其中包含合法文件 ‘valid.c’。
    \item \textbf{错误输入示例}：路径不存在、路径下无目标文件、文件内括号匹配错误等。
\end{itemize}

\section{概要设计}

\subsection{问题解决思路}
括号匹配是典型的栈应用问题。程序采用以下步骤：
\begin{enumerate}
    \item 用户输入目录路径，程序使用 Windows API 函数 ‘FindFirstFile’ / ‘FindNextFile’遍历目录，筛选出扩展名为 ‘.c’、‘.cpp’、‘.java’ 的文件，将完整路径存入动态数组。
    \item 对每个文件依次处理：
        \begin{itemize}
            \item 逐行读取文件内容。
            \item 遇到左括号 ‘(’、‘[’、‘{’时，将其字符和所在行号压入栈。
            \item 遇到右括号 ‘)’、‘]’、‘}’ 时，检查栈是否为空：若空则报告“未匹配的右括号”；否则检查栈顶左括号是否与当前右括号匹配：若匹配则弹出栈顶，否则报告“嵌套顺序错误”。
            \item 若文件读取完毕栈非空，则报告栈底左括号所在行号为“未匹配的左括号”。
            \item 一旦发现任何错误，立即停止当前文件的检查，输出错误信息后继续处理下一个文件。
        \end{itemize}
    \item 处理完成后释放动态数组内存。
\end{enumerate}

\subsection{数据结构定义}
\begin{itemize}
    \item \textbf{节点结构体 Node}：存储括号字符及其所在行号。
\begin{lstlisting}
typedef struct {
    char ch;      // 括号字符
    int line;     // 所在行号
} Node;
\end{lstlisting}
    \item \textbf{栈结构体 Stack}：使用定长数组存储节点，‘top’ 为栈顶指针。
\begin{lstlisting}
typedef struct {
    Node data[MAX_DEPTH];   // MAX_DEPTH = 10000
    int top;                // 栈顶索引，初始 -1
} Stack;
\end{lstlisting}
\end{itemize}

\subsection{主程序流程}
\begin{enumerate}
    \item 提示用户输入路径，读取并去除换行符。
    \item 调用 ‘scan\_files’扫描目录，获取目标文件路径列表。
    \item 若文件列表为空，输出提示并退出。
    \item 对列表中的每个文件：
        \begin{enumerate}
            \item 打开文件，若失败则跳过并提示。
            \item 初始化栈 ‘braket’，错误标志 ‘error = 0’，行号计数器 ‘num = 0’。
            \item 逐行读取：
                \begin{itemize}
                    \item 遍历行内每个字符，判断是否为括号。
                    \item 左括号：压栈（若栈满则报告嵌套过深）。
                    \item 右括号：根据栈的状态进行匹配或报错。
                    \item 发现错误则 ‘break‘ 退出内层循环。
                \end{itemize}
            \item 若文件正常读完且栈为空，输出合法信息；若栈非空，输出栈底左括号对应的未匹配错误。
            \item 关闭文件。
        \end{enumerate}
    \item 释放文件路径数组，程序结束。
\end{enumerate}

\subsection{模块层次关系}
\begin{verbatim}
main
├── scan_files          // 扫描目录，构建文件路径列表
│   └── IsTargetFile    // 判断文件扩展名是否为目标类型
├── 对每个文件循环
│   ├── fopen / fclose
│   ├── push / pop / isBra / isKet / match
│   └── ErrorPrint      // 根据错误类型输出信息
└── free_filepaths      // 释放动态数组
\end{verbatim}

\section{详细设计}

\subsection{数据类型定义}
已在前文给出 ‘Node’ 和 ‘Stack’ 结构体。此外定义了以下常量和辅助函数：
\begin{lstlisting}
#define MAX_DEPTH 10000      // 最大嵌套深度
#define MAX_LINE_LEN 4096    // 每行最大长度
#define MAX_FILE 1000        // 最大文件数量（实际动态分配）
\end{lstlisting}

\subsection{核心算法伪代码}

\subsubsection{栈操作}
\begin{lstlisting}
function push(Stack* s, char ch, int line):
    if s->top + 1 < MAX_DEPTH:
        s->top++
        s->data[s->top].ch = ch
        s->data[s->top].line = line
        return 0
    else:
        return -1   // 栈满

function pop(Stack* s):
    if s->top >= 0:
        s->top--

function isBra(char ch):
    return ch == '(' or ch == '[' or ch == '{'

function isKet(char ch):
    return ch == ')' or ch == ']' or ch == '}'

function match(char left, char right):
    return (left=='(' and right==')') or
           (left=='[' and right==']') or
           (left=='{' and right=='}')
\end{lstlisting}

\subsubsection{目录扫描（Windows 版本，借助AI完成）}
\begin{lstlisting}
function scan_files(dir_path, &filepaths, &count):
    search_path = dir_path + "\\*"
    hFind = FindFirstFile(search_path, &find_data)
    if hFind == INVALID_HANDLE_VALUE: return -1
    capacity = 10
    filepaths = malloc(capacity * sizeof(char*))
    count = 0
    do:
        if find_data.cFileName is "." or "..": continue
        if not a directory and IsTargetFile(find_data.cFileName):
            full_path = dir_path + "\\" + find_data.cFileName
            path_copy = _strdup(full_path)
            // 动态扩容
            if count >= capacity:
                capacity *= 2
                realloc filepaths
            filepaths[count++] = path_copy
    while FindNextFile(hFind, &find_data) != 0
    FindClose(hFind)
    return 0
\end{lstlisting}

\subsubsection{括号匹配主循环}
\begin{lstlisting}
for each file in filepaths:
    fp = fopen(file, "r")
    if fp == NULL: continue
    Stack braket; braket.top = -1
    error = 0
    lineNum = 0
    while fgets(curline, MAX_LINE_LEN, fp):
        lineNum++
        for i = 0 to strlen(curline)-1:
            c = curline[i]
            if isBra(c):
                if push(&braket, c, lineNum) != 0:
                    print "括号嵌套过深"
                    error = 1; break
            else if isKet(c):
                if braket.top == -1:
                    ErrorPrint(3, lineNum, c)   // 多余右括号
                    error = 1; break
                else if match(braket.data[braket.top].ch, c):
                    pop(&braket)
                else:
                    ErrorPrint(2, lineNum, c)   // 嵌套顺序错误
                    error = 1; break
        if error: break
    fclose(fp)
    if error == 0 and braket.top == -1:
        print "文件[%s]括号匹配合法"
    else if braket.top >= 0:
        // 输出栈底左括号（最早出现）
        ErrorPrint(4, braket.data[0].line, braket.data[0].ch)
\end{lstlisting}

\subsection{函数调用关系图}
\begin{verbatim}
main
├── scan_files
│   ├── IsTargetFile
│   └── _strdup / malloc / realloc
├── 文件处理循环
│   ├── push
│   ├── pop
│   ├── isBra / isKet
│   ├── match
│   └── ErrorPrint
└── free_filepaths
\end{verbatim}

\section{调试分析报告}

\subsection{调试过程中遇到的问题及解决方法}
\begin{enumerate}
    \item \textbf{问题1：字符串遍历越界} \\
    初始版本中使用 ‘while(*p)’ 配合 ‘p[i]’ 且未移动 ‘p’，导致循环条件永远为真，访问越界。 \\
    \textbf{解决方法}：改用 ‘for (int i = 0; curline[i]; i++)’ 直接遍历字符数组，避免指针混淆。
    
    \item \textbf{问题2：每行结束后错误清空栈} \\
    最初代码在每行读取完毕后，若栈非空就输出“未匹配的左括号”并清空栈，导致跨行括号匹配失败。 \\
    \textbf{解决方法}：删除该段逻辑，将栈剩余括号的检查移至文件全部读取之后。
    
    \item \textbf{问题3：遇到错误未立即停止当前文件} \\
    发现括号错误后仅用 ‘break’ 退出内层循环，但外层 ‘fgets’ 继续执行，导致同一文件输出多个错误。 \\
    \textbf{解决方法}：在 ‘break’ 后增加 ‘if(error) break;’ 跳出文件读取循环。
    
    \item \textbf{问题4：非目标文件导致程序退出} \\
    最初遇到非 ‘.c/.cpp/.java’ 文件时会打印错误并 ‘return -1’，不符合“批量处理”要求。 \\
    \textbf{解决方法}：在 ‘scan\_ files’ 阶段直接过滤非目标文件，主循环不再检查类型；同时 ‘fopen‘ 失败时使用 ‘continue’ 跳过。
\end{enumerate}

\subsection{算法时空复杂度分析}
\begin{itemize}
    \item \textbf{时间复杂度}：设共有 \( f \) 个目标文件，每个文件有 \( l \) 行，平均每行 \( len \) 个字符。程序对每个字符进行常数次判断和栈操作，故总时间复杂度为 \( O(f \times l \times len) \)，即 $ O($文件总字符数$) $。实际运行效率很高。
    \item \textbf{空间复杂度}：栈最大深度为常数 MAX\_ DEPTH$=10000 $，文件路径数组大小与文件数量 \( f \) 成正比，故空间复杂度为 \( O(f) \)。栈内每个节点存储一个字符和一个整数，占用常数空间。
\end{itemize}

\subsection{改进设想}
\begin{itemize}
    \item 支持递归遍历子目录（当前仅处理指定目录下的文件）。
    \item 增加对注释和字符串字面量的忽略，避免注释内的括号被误判（实验假设无注释，故暂未实现）。
    \item 提供命令行参数方式直接传入路径，无需交互输入。
\end{itemize}

\subsection{经验与体会}
通过本次实验，我深入理解了栈在括号匹配中的经典应用，掌握了动态内存管理的技巧。调试过程中对指针和循环边界的细致检查，提升了代码排错能力。同时，严格按照实验要求输出首个错误位置，锻炼了需求分析和规范编码的能力。

\section{用户使用说明}
\begin{enumerate}
    \item 本程序运行于 Windows 操作系统（使用了 ‘<windows.h>’），需使用支持 C99 的编译器（如 MinGW、Visual Studio）编译。
    \item 编译命令示例（MinGW）：\texttt{gcc chap3.exp.c -o bracket\_checker.exe}
    \item 运行程序：点击“编译运行”图标或在命令行中执行。
    \item 程序提示“请输入一个路径：”时，输入待检查的目录路径（例如 \texttt{D:/mycode}），注意路径中不要包含中文（避免编码问题），按回车确认。
    \item 程序将自动扫描该目录下的所有 \texttt{.c}、\texttt{.cpp}、\texttt{.java} 文件，忽略无关文件，并依次输出每个文件的校验结果。
    \item 若目录不存在或无合法文件，程序会给出相应提示并退出。
    \item 对于每个非法文件，仅输出首个错误所在行号和错误类型，然后继续检查下一个文件。
    \item 检查完毕后程序自动结束。
\end{enumerate}

\textbf{注意事项}：
\begin{itemize}
    \item 程序假设代码文件中没有注释，若实际文件包含注释，注释内的括号也会被检查，可能导致误报。
    \item 路径中不要使用反斜杠转义，直接输入如 \texttt{C:/test} 即可。
    \item 若文件编码不是系统默认 ANSI（如 UTF-8 without BOM），可能导致中文乱码，建议将源文件转为 GBK 编码或修改终端代码页。
\end{itemize}

\section{测试结果}

\subsection{测试环境}
\begin{itemize}
    \item 操作系统：Windows 11
    \item 编译器：MinGW-W64 GCC 8.1.0
    \item 编译选项：\texttt{gcc chap3.exp.c -o test.exe}
\end{itemize}

\subsection{测试用例1：全部合法文件}
目录 ‘D:/test/valid’ 下包含三个文件：
\begin{itemize}
    \item \texttt{valid1.c}：内容为 \texttt{int main() \{ printf("Hello"); \}}
    \item \texttt{valid2.cpp}：内容为 \texttt{void func() \{ if (a[b]) \{ \} \}}
    \item \texttt{valid3.java}：内容为 \texttt{class A \{ void m() \{ [] \} \}}
\end{itemize}
\textbf{输出结果}：
\begin{verbatim}
文件[D:\test\valid\valid1.c]括号匹配合法
文件[D:\test\valid\valid2.cpp]括号匹配合法
文件[D:\test\valid\valid3.java]括号匹配合法
\end{verbatim}

\subsection{测试用例2：包含括号错误}
目录 ‘D:/test/invalid’ 下：
\begin{itemize}
    \item \texttt{error1.c} 内容：
\begin{lstlisting}
1: #include <stdio.h>
2: void main() {
3:     if (a == 1] {   // 方括号与圆括号不匹配
4:         printf("error");
5:     }
6: }
\end{lstlisting}
    \item \texttt{error2.cpp} 内容：
\begin{lstlisting}
1: int main() {
2:     (
3:     return 0;
4: }
\end{lstlisting}
\end{itemize}
\textbf{输出结果}：
\begin{verbatim}
第3行：方括号嵌套顺序错误
第2行：未匹配的左圆括号
\end{verbatim}

\subsection{测试用例3：边界情况}
\begin{itemize}
    \item 空文件 \texttt{empty.c}：输出 \texttt{文件[empty.c]括号匹配合法}
    \item 只有左括号的文件 \texttt{onlyleft.c}：内容为 \texttt{\{ ( [ }，输出 \texttt{第1行：未匹配的左花括号}
    \item 只有右括号的文件 \texttt{onlyright.c}：内容为 \texttt{ ) ] }，输出 \texttt{第1行：未匹配的右圆括号}
\end{itemize}

\subsection{测试用例4：路径不存在}
输入路径 \texttt{C:/notexist}：
\begin{verbatim}
无法打开目录: C:/notexist
扫描目录失败
\end{verbatim}

\subsection{测试用例5：目录下无目标文件}
输入路径 ‘D:/empty’（该目录下只有 ‘.txt’ 文件）：
\begin{verbatim}
目录下没有文件
\end{verbatim}

所有测试均通过，程序行为符合实验要求。

\end{document}